\chapter{Fontes de dados}\label{anx:data}

Todos os conjuntos de dados provêm do catálogo público da plataforma de computação, acessados
em resolução nativa e reescalados para unidades físicas quando o provedor armazena bandas
empacotadas em inteiros; cada fonte foi verificada contra o catálogo ativo no momento do
processamento.

\begin{table}[H]
\centering
\caption{Fontes de dados biofísicos e agroclimáticos.}
\label{tab:src-biophys}
\begin{tabular}{@{}p{2.7cm}p{3.6cm}p{2cm}p{5.5cm}@{}}
\toprule
Tema & Conjunto de dados (provedor) & Resolução nativa & Variáveis utilizadas \\
\midrule
Clima e balanço hídrico & TerraClimate, mensal & $\approx$4~km & Precipitação,
  evapotranspiração potencial e real, déficit hídrico climático, umidade do solo, déficit de
  pressão de vapor, radiação solar, temperatura máxima/mínima \\
Terreno & Modelo digital de elevação SRTM & 30~m & Elevação, declive, orientação, posição
  topográfica \\
Hidrografia & Acumulação de fluxo HydroSHEDS & 15 segundos de arco & Índice de umidade
  topográfica, densidade de drenagem \\
Solo (0--30~cm) & Propriedades de solo OpenLandMap & 250~m & Argila, areia, carbono orgânico,
  pH (H$_2$O), densidade aparente, teor de água disponível para as plantas \\
Água de superfície & JRC Global Surface Water & 30~m & Ocorrência e sazonalidade
  $\rightarrow$ distância à água permanente \\
Fenologia da vegetação & MODIS NDVI, composições de 16 dias & 250~m & Parâmetros harmônicos
  (média, amplitude, época de pico, integral) \\
Acessibilidade ao mercado & Global Accessibility to Cities (2015) & $\approx$1~km & Tempo de
  viagem à cidade mais próxima (\emph{proxy} de acessibilidade ao mercado) \\
Cobertura genérica da terra & MapBiomas, coleção~10 (2024) & 30~m & Máscara e frações por classe
  (apenas máscara/contexto) \\
Áreas protegidas & World Database on Protected Areas (WDPA) & vetor & Distância a reservas
  $\rightarrow$ conectividade de conservação (camada de governança, não de cobertura da
  terra) \\
Biomassa acima do solo & Densidade global de carbono em biomassa acima do solo (Spawn
  \emph{et al.}) & $\approx$300~m & Valor do serviço de armazenamento de carbono
  (conservação) \\
Clima futuro & NASA NEX-GDDP-CMIP6, diário com redução de escala & $\approx$25~km &
  Precipitação, temperatura máxima e mínima $\rightarrow$ fatores de mudança
  (\emph{delta-change}) \\
\bottomrule
\end{tabular}
\end{table}

As bandas de solo são a média de 0--30~cm dos horizontes de 0, 10 e 30~cm. Conversões de
unidade verificadas: carbono orgânico do solo $\times 5$ (média territorial $\approx$
13~g~kg$^{-1}$), pH $\times 0{,}1$ ($\approx$4,7--7,3) e densidade aparente $\times 10$
($\approx$1259~kg~m$^{-3}$); argila, areia e teor de água não requerem reescalonamento.

\begin{table}[H]
\centering
\caption{Fontes de dados de uso atual e validação.}
\label{tab:src-realized}
\begin{tabular}{@{}p{5cm}p{2.7cm}p{6cm}@{}}
\toprule
Fonte (provedor) & Resolução nativa & Finalidade \\
\midrule
Malha Municipal Digital do IBGE (2025) & $\approx$111~m (simplificado, em memória) & Agregação e
  relato municipais (247 municípios; sem \emph{upload}) \\
Uso e cobertura da terra anual MapBiomas Brasil & 30~m & Máscaras, análise potencial
  \emph{versus} realizado, caracterização descritiva da composição \\
Produtividade primária líquida (anual) e bruta (8 dias) do MODIS & 500~m & Validação da
  viabilidade (substituto de rendimento) \\
\bottomrule
\end{tabular}
\end{table}

O uso atual da terra e o \emph{proxy} de produtividade são usados estritamente para
mascaramento, validação e caracterização descritiva; nunca entram nos cálculos de viabilidade
ou de agrupamento.

\begin{table}[H]
\centering
\caption{Especificação do conjunto de modelos CMIP6.}
\label{tab:cmip6-spec}
\begin{tabular}{@{}p{4.5cm}p{9.5cm}@{}}
\toprule
Parâmetro & Valor \\
\midrule
Modelos de circulação geral & ACCESS-CM2, EC-Earth3, GFDL-ESM4, MPI-ESM1-2-HR, MRI-ESM2-0 \\
Cenários & SSP2-4.5, SSP5-8.5 \\
Janelas futuras & 2031--2050 e 2051--2070 (inclusive) \\
Linha de base histórica dos modelos & 1991--2014 (o experimento histórico com redução de
  escala termina em 2014) \\
Variáveis & Precipitação (fator de mudança multiplicativo); temperatura máxima e mínima
  (fator de mudança aditivo, em kelvin) \\
\bottomrule
\end{tabular}
\end{table}

\chapter{Especificação completa de viabilidade e pesos AHP}\label{anx:spec}

A Tabela~\ref{tab:spec-longtable} reproduz a especificação completa por segmento subjacente à
\S\ref{sec:suitability}. O declive está em graus. Os limiares são valores de primeira
aproximação ancorados na literatura e reconciliados com a faixa realizada de cada variável
sobre GO~e~DF pela recalibração empírica; os pesos, por sua vez, são reconciliados com o poder
discriminante amostral de cada fator (\S\ref{sec:ahp}), de modo que um fator quase constante em
GO~e~DF --- por mais importante que a literatura o considere --- carrega um peso baixo.

\begin{longtable}{@{}p{2.6cm}p{3.6cm}p{1cm}p{1.7cm}p{3.3cm}@{}}
\caption{Fatores, pesos, formas de pertinência e limiares por segmento.}\label{tab:spec-longtable} \\
\toprule
Segmento (máscara) & Fator & Peso & Forma & Limiares \\
\midrule
\endfirsthead
\toprule
Segmento (máscara) & Fator & Peso & Forma & Limiares \\
\midrule
\endhead
% GENERATED by tools/gen_thesis_tables.py -- do not edit by hand.
% Edit the source CSV and regenerate; see CLAUDE.md §9.
\textbf{Soja} (disponível) & pH do solo & 0,38 & intervalar & 4,5, 5,5, 6,8, 7,5 \\
 & Declive ($^\circ$) & 0,37 & decrescente & 4, 8 \\
 & Teor de argila (\%) & 0,24 & intervalar & 20, 32, 45, 62 \\
 & Temp. do trimestre mais quente ($^\circ$C) & porta & intervalar & 18, 22, 30, 34 \\
 & Precipitação anual (mm) & porta & intervalar & 800, 1100, 1800, 2200 \\
 & Índice de aridez (P/ETP) & porta & intervalar & 0,4, 0,7, 1,3, 1,6 \\
\textbf{Cana-de-açúcar} (disponível) & Teor de argila (\%) & 0,38 & intervalar & 22, 32, 52, 66 \\
 & Teor de água disponível para as plantas & 0,32 & crescente & 22,1, 33,6 \\
 & Declive ($^\circ$) & 0,30 & decrescente & 5, 10 \\
 & Temp. do trimestre mais quente ($^\circ$C) & porta & intervalar & 20, 23, 32, 36 \\
 & Precipitação anual (mm) & porta & intervalar & 900, 1300, 2000, 2600 \\
 & Meses secos (n) & porta & intervalar & 3, 5, 5, 9 \\
\textbf{Outras culturas anuais} (disponível) & Carbono orgânico do solo (g/kg) & 0,54 & crescente & 6, 20 \\
 & pH do solo & 0,16 & intervalar & 4,5, 5,5, 6,8, 7,5 \\
 & Graus-dia de crescimento & 0,14 & intervalar & 2500, 3500, 5500, 7000 \\
 & Declive ($^\circ$) & 0,12 & decrescente & 4, 8 \\
 & CV da precipitação & 0,04 & decrescente & 0,6, 1,3 \\
 & Precipitação anual (mm) & porta & intervalar & 700, 1000, 1700, 2200 \\
\textbf{Piscicultura} (disponível) & Distância à água sazonal (m) & 0,48 & decrescente & 750, 6000 \\
 & Temp. do trimestre mais quente ($^\circ$C) & 0,22 & intervalar & 22, 26, 32, 36 \\
 & Teor de argila (\%) & 0,14 & crescente & 15, 35 \\
 & Declive ($^\circ$) & 0,11 & decrescente & 2,86, 5,71 \\
 & Densidade de drenagem & 0,05 & crescente & 0, 0,09 \\
\textbf{Pecuária} (apenas exclusão) & Umidade do solo (mm) & 0,70 & crescente & 55, 120 \\
 & Declive ($^\circ$) & 0,17 & decrescente & 6, 15 \\
 & Carbono orgânico do solo (g/kg) & 0,13 & crescente & 8, 18 \\
 & Precipitação anual (mm) & porta & intervalar & 700, 1000, 2000, 2600 \\
 & Meses secos (n) & porta & decrescente & 5, 9 \\
\textbf{Conservação} (nenhuma; compensatória) & Distância a área protegida (m) & 0,35 & decrescente & 5000, 40000 \\
 & Carbono em biomassa acima do solo (Mg~C/ha) & 0,17 & crescente & 10, 20 \\
 & Rugosidade do terreno (m) & 0,17 & crescente & 5, 30 \\
 & Índice de umidade topográfica & 0,07 & crescente & 9,44, 11,56 \\
 & Fração de vegetação nativa & 0,07 & crescente & 0,25, 0,85 \\
 & Índice de aridez (P/ETP) $\dagger$ & 0,06 & crescente & 0,95, 1,25 \\
 & Temp. do trimestre mais quente ($^\circ$C) $\dagger$ & 0,05 & decrescente & 23, 31 \\
 & Densidade de drenagem & 0,03 & crescente & 0,01, 0,08 \\
 & Declive ($^\circ$) & 0,02 & crescente & 2,86, 11,31 \\
\textbf{Geração fotovoltaica} (disponível) & Radiação solar (W/m$^2$) & 0,40 & crescente & 186, 204 \\
 & Temp. do trimestre mais quente ($^\circ$C) & 0,18 & decrescente & 24, 31 \\
 & Log do tempo de viagem à cidade & 0,14 & decrescente & 3,5, 5,4 \\
 & Índice de claridade (cobertura de nuvens) & 0,14 & crescente & 0,47, 0,5 \\
 & Declive ($^\circ$) & 0,08 & decrescente & 2,86, 8,53 \\
 & Componente norte da orientação & 0,06 & crescente & -0,3, 0,5 \\
\bottomrule
   % the fragment ends with \bottomrule
\end{longtable}

A conservação ($\dagger$) carrega dois fatores climáticos \textbf{móveis} (aridez e
temperatura do trimestre mais quente), que a tornam responsiva ao CMIP6; é agregada pela
média aritmética ponderada \textbf{compensatória} (\S\ref{sec:aggregation}), e sua fração de
vegetação nativa (peso 0,07) é a única entrada de cobertura da terra, rebaixada, sancionada
pela salvaguarda contra circularidade. Os outros seis segmentos usam a média geométrica de
fator limitante.

%% TODO: é mais facil remover estas variaveis do que aumentar a complexidade da analise
Os fatores marcados como \emph{porta} não entram na matriz AHP e não recebem peso: são os
fatores cuja pertinência é constante em GO~e~DF (a faixa observada cai inteiramente dentro do
seu platô ótimo), aplicados multiplicativamente com limiares ancorados de modo que o valor
presente dê $\mu = 1$ (\S\ref{sec:ahp}). Por isso a soja e a cana-de-açúcar têm três fatores
ponderados, e não seis, e a pecuária três, e não cinco: os demais foram promovidos a porta,
não removidos.

Os pesos de cada segmento são as prioridades de uma matriz Saaty de comparação par a par
documentada; toda matriz é folgadamente consistente, sendo 0,0462 a maior razão de
consistência entre os sete segmentos (Tabela~\ref{tab:ahp-cr}), menos da metade do limiar de
aceitação de 0,10.

\begin{table}[H]
\centering
\caption{Razão de consistência AHP por segmento.}
\label{tab:ahp-cr}
% GENERATED by tools/gen_thesis_tables.py -- do not edit by hand.
% Edit the source CSV and regenerate; see CLAUDE.md §9.
\begin{tabular}{@{}lrr@{}}
\toprule
Segmento & Fatores ($n$) & Razão de consistência \\
\midrule
Soja & 6 & 0,0133 \\
Cana-de-açúcar & 6 & 0,0335 \\
Outras culturas anuais & 6 & 0,0054 \\
Piscicultura & 5 & 0,0168 \\
Pecuária & 5 & 0,0067 \\
Conservação & 9 & 0,0245 \\
Geração fotovoltaica & 5 & 0,0205 \\
\bottomrule
\end{tabular}

\end{table}

\chapter{Derivação dos atributos}\label{anx:derivations}

\textbf{Clima e balanço hídrico (14 variáveis).} Uma climatologia de 12 meses é computada
para cada banda climática sobre o registro de 1991--2020 (a média climatológica de cada mês
do calendário); os agregados anuais são a soma das doze médias mensais. As variáveis
derivadas do balanço hídrico são a precipitação anual, a evapotranspiração potencial anual, o
déficit hídrico climático, a evapotranspiração real e o índice de aridez (precipitação anual
$\div$ evapotranspiração potencial anual). As variáveis de sazonalidade são a sazonalidade da
precipitação (o coeficiente de variação das doze médias mensais), a duração da estação seca
(meses com precipitação média abaixo de 60~mm) e a precipitação da estação chuvosa (a soma
sobre os meses mais úmidos que a média mensal geral). Completam o conjunto a umidade média do
solo, o déficit de pressão de vapor e a radiação solar médios, os graus-dia de crescimento, e
a temperatura média do trimestre mais quente e do mais frio (o máximo e o mínimo das doze
médias móveis de três meses). Os graus-dia de
crescimento são acumulados a partir da climatologia mensal de temperatura média com uma
temperatura base $T_{\text{base}} = 10\,^\circ\mathrm{C}$,
\begin{equation*}
  \mathrm{GDD} = \sum_{m=1}^{12} \max\!\left(\bar{T}_m - T_{\text{base}},\, 0\right)\, N_m,
\end{equation*}
onde $\bar{T}_m$ é a temperatura média do mês $m$ e $N_m = 30{,}4$ dias, a duração média do mês. É o
Método~1 de \citet[p.~292]{mcmaster_growing_1997} aplicado a médias mensais, o que tende a
subestimar o GDD em meses próximos da temperatura base. Toda a
derivação climática é encapsulada em uma única transformação determinística reutilizada sem
alteração para o clima futuro do CMIP6 (\S\ref{sec:cmip6}), o que é o que torna a mudança
projetada um sinal puramente climático.

\textbf{Terreno (6 variáveis).} A elevação e o declive vêm do MDE de 30~m; a orientação é
decomposta em uma componente norte ($\cos$) e leste ($\sin$), para que a variável circular
seja utilizável em aritmética de pixels. O índice de posição topográfica é a diferença entre
a elevação $z_0$ de uma célula e a elevação média de sua vizinhança circular,
$\mathrm{TPI} = z_0 - \overline{z}_{\text{focal}}$ (positivo em cristas, negativo em vales). O
índice de umidade topográfica é $\mathrm{TWI} = \ln(\alpha / \tan\beta)$
\citep[p.~48, eq.~8]{beven_physically_1979}, com $\beta$ o declive local e $\alpha$ a área de
contribuição específica, aproximada por $(n+1)\,\Delta x$, em que $n$ é a acumulação de fluxo do
HydroSHEDS em número de células e $\Delta x \approx 464$~m a largura da sua célula de 15 segundos de
arco.

\textbf{Solo (6 variáveis).} A média de 0--30~cm dos horizontes de 0, 10 e 30~cm para argila,
areia, carbono orgânico, pH (H$_2$O), densidade aparente e teor de água disponível para as
plantas, cada um reescalado para unidades físicas pelo seu fator de conversão verificado
(Anexo~\ref{anx:data}).

\textbf{Água (3 variáveis).} Uma máscara de água permanente (ocorrência de água superficial
$>$ 80\%) produz uma superfície de distância euclidiana à água na grade de 250~m; a
sazonalidade da água de superfície é carregada do produto de origem; a densidade de drenagem
é a fração de células de canal (acumulação de fluxo acima de um limiar fixo) em uma
vizinhança circular, com faixa empírica de $\approx$0--0,15 (percentil 90 $\approx$0,095) que
governa seus limiares de pertinência.

\textbf{Fenologia (4 variáveis).} Os parâmetros fenológicos são extraídos por uma regressão
harmônica (de Fourier) do ciclo sazonal de NDVI, ajustada à climatologia de NDVI de 12 meses
em vez das $\sim$460 composições de 16 dias ($\approx$38$\times$ mais barato para as mesmas
métricas). O modelo mantém um termo constante mais harmônicos anual e semianual,
\begin{equation*}
  \mathrm{NDVI}(t) = \beta_0 + \beta_1 \cos(2\pi t) + \gamma_1 \sin(2\pi t)
   + \beta_2 \cos(4\pi t) + \gamma_2 \sin(4\pi t) + \varepsilon(t),
\end{equation*}
com $t$ a fração do ano, fornecendo o NDVI médio ($\beta_0$), a amplitude sazonal
$A = \sqrt{\beta_1^2 + \gamma_1^2}$, a época de pico (a partir da fase do harmônico anual) e
uma integral anual (NDVI médio $\times$ 365).

\textbf{Acessibilidade ao mercado (1 variável).} A superfície de tempo de viagem às cidades é
transformada em logaritmo como $\ln(t+1)$ e reamostrada para a grade.

\textbf{Máscara de cobertura da terra (apenas contexto).} O MapBiomas, a 30~m, é
agregado a 250~m como uma classe majoritária mais frações por classe, produzindo uma máscara de
áreas \emph{excluídas} (urbano, água permanente, área úmida herbácea), uma máscara de áreas
\emph{disponíveis} (lavoura, campo, vegetação arbustiva, solo exposto) e frações de
árvore/lavoura/campo para contexto.

\chapter{Tabelas de resultados suplementares}\label{anx:tables}

\begin{table}[H]
\centering
\caption{Médias biofísicas discriminantes por zona.}
\label{tab:zone-means}
% GENERATED by tools/gen_thesis_tables.py -- do not edit by hand.
% Edit the source CSV and regenerate; see CLAUDE.md §9.
\begin{tabular}{@{}lrrrr@{}}
\toprule
Variável & Z1 & Z2 & Z3 & Z4 \\
\midrule
Tamanho da amostra (pontos) & 10.163 & 2.768 & 6.089 & 7.800 \\
Parcela da área de interesse & 46,4\% & 10,4\% & 15,2\% & 28,1\% \\
Elevação (m) & 784 & 733 & 483 & 445 \\
Declive ($^\circ$) & 4,5 & 16,9 & 3,5 & 3,3 \\
Índice de umidade topográfica & 10,2 & 8,8 & 11,4 & 10,6 \\
Temp. do trimestre mais quente ($^\circ$C) & 24,5 & 25,5 & 25,6 & 26,6 \\
Índice de aridez (P/ETP) & 1,15 & 1,10 & 1,13 & 1,05 \\
Umidade do solo (mm) & 86 & 110 & 93 & 109 \\
Teor de argila (\%) & 37,0 & 34,4 & 29,4 & 26,4 \\
Carbono orgânico do solo (g/kg) & 15,0 & 16,0 & 13,5 & 11,1 \\
Densidade de drenagem (fração) & 0,006 & 0,009 & 0,107 & 0,012 \\
Distância à água (m) & 16.181 & 19.624 & 8.510 & 15.620 \\
\bottomrule
\end{tabular}

\end{table}

\begin{table}[H]
\centering
\caption{Mudança média projetada $\Delta S$ por zona para as três culturas sensíveis ao clima,
nos quatro cenário--janela.}
\label{tab:zone-shift}
% GENERATED by tools/gen_thesis_tables.py -- do not edit by hand.
% Edit the source CSV and regenerate; see CLAUDE.md §9.
\begin{tabular}{@{}llrrrr@{}}
\toprule
Segmento & Cenário--janela & Z1 & Z2 & Z3 & Z4 \\
\midrule
Soja & SSP2-4.5 2031--2050 & $+$0,000 & $+$0,000 & $-$0,000 & $-$0,000 \\
Soja & SSP2-4.5 2051--2070 & $-$0,000 & $-$0,000 & $-$0,000 & $-$0,000 \\
Soja & SSP5-8.5 2031--2050 & $-$0,000 & $+$0,000 & $-$0,000 & $-$0,000 \\
Soja & SSP5-8.5 2051--2070 & $-$0,000 & $-$0,000 & $-$0,001 & $-$0,004 \\
\midrule
Cana-de-açúcar & SSP2-4.5 2031--2050 & $+$0,000 & $-$0,000 & $-$0,000 & $-$0,001 \\
Cana-de-açúcar & SSP2-4.5 2051--2070 & $+$0,000 & $-$0,000 & $-$0,000 & $-$0,001 \\
Cana-de-açúcar & SSP5-8.5 2031--2050 & $-$0,001 & $-$0,000 & $-$0,000 & $-$0,001 \\
Cana-de-açúcar & SSP5-8.5 2051--2070 & $-$0,000 & $-$0,000 & $-$0,000 & $-$0,002 \\
\midrule
Outras culturas anuais & SSP2-4.5 2031--2050 & $-$0,004 & $-$0,003 & $-$0,012 & $-$0,027 \\
Outras culturas anuais & SSP2-4.5 2051--2070 & $-$0,012 & $-$0,006 & $-$0,025 & $-$0,048 \\
Outras culturas anuais & SSP5-8.5 2031--2050 & $-$0,008 & $-$0,005 & $-$0,018 & $-$0,038 \\
Outras culturas anuais & SSP5-8.5 2051--2070 & $-$0,036 & $-$0,012 & $-$0,053 & $-$0,093 \\
\bottomrule
\end{tabular}

\end{table}

%% TODO: melhorar cabeçalho dessa tabela. Colocar cenario SSP em celulas mescladas na primeira linha e
%% TODO: janela temporal na segunda linha, de forma a economizar espaço horizontal

\begin{table}[H]
\centering
\caption{Parcela da terra avaliada de cada segmento rebaixada em pelo menos uma classe da
FAO, presente $\rightarrow$ futuro (promoções desprezíveis e omitidas).}
\label{tab:downgrade-share}
\small
% GENERATED by tools/gen_thesis_tables.py -- do not edit by hand.
% Edit the source CSV and regenerate; see CLAUDE.md §9.
\begin{tabular}{@{}lrrrr@{}}
\toprule
Segmento & SSP2-4.5 2031--2050 & SSP2-4.5 2051--2070 & SSP5-8.5 2031--2050 & SSP5-8.5 2051--2070 \\
\midrule
Soja & 0,0\% & 0,1\% & 0,0\% & 1,0\% \\
Cana-de-açúcar & 0,1\% & 0,1\% & 0,3\% & 0,2\% \\
Outras culturas anuais & 4,8\% & 9,5\% & 7,2\% & 21,8\% \\
Piscicultura & 0,0\% & 0,0\% & 0,0\% & 0,0\% \\
Pecuária & 0,0\% & 0,0\% & 0,0\% & 0,5\% \\
Conservação & 14,5\% & 16,8\% & 17,4\% & 25,6\% \\
Geração fotovoltaica & 11,5\% & 18,3\% & 15,3\% & 33,3\% \\
\bottomrule
\end{tabular}

\end{table}

\begin{table}[H]
\centering
\caption{Matriz de transição de classes da FAO, presente $\rightarrow$ futuro, para as
outras culturas anuais sob SSP5-8.5, 2051--2070 (área como percentual da terra avaliada do
segmento).}
\label{tab:transition-othercrops}
% GENERATED by tools/gen_thesis_tables.py -- do not edit by hand.
% Edit the source CSV and regenerate; see CLAUDE.md §9.
\begin{tabular}{@{}lrrrrr@{}}
\toprule
Presente $\downarrow$ / futuro $\rightarrow$ & N & S3 & S2 & S1 & Total presente \\
\midrule
N & 36,5\% & 0,0\% & 0,0\% & 0,0\% & \textbf{36,5\%} \\
S3 & 2,2\% & 24,7\% & 0,0\% & 0,0\% & \textbf{26,9\%} \\
S2 & 0,1\% & 6,3\% & 19,6\% & 0,0\% & \textbf{25,9\%} \\
S1 & 0,0\% & 0,0\% & 3,0\% & 7,7\% & \textbf{10,7\%} \\
\bottomrule
\end{tabular}

\end{table}

\begin{table}[H]
\centering
\caption{Composição do uso atual da terra dentro de cada zona biofísica (parcela da área da
zona).}
\label{tab:zone-realized}
% GENERATED by tools/gen_thesis_tables.py -- do not edit by hand.
% Edit the source CSV and regenerate; see CLAUDE.md §9.
\begin{tabular}{@{}lrrrr@{}}
\toprule
Uso atual & Z1 & Z2 & Z3 & Z4 \\
\midrule
Soja & 23,1\% & 0,1\% & 7,9\% & 5,4\% \\
Cana-de-açúcar & 3,5\% & 0,0\% & 3,1\% & 0,6\% \\
Outras culturas anuais & 0,8\% & 0,0\% & 1,0\% & 1,3\% \\
Pastagem & 44,0\% & 29,7\% & 58,1\% & 66,5\% \\
Vegetação nativa & 25,0\% & 69,7\% & 22,7\% & 23,7\% \\
\bottomrule
\end{tabular}

\end{table}

\begin{table}[H]
\centering
\caption{Correlação de Spearman da viabilidade média municipal com o \emph{proxy} de NPP
anual do MODIS (terra manejada; $n = 247$ municípios).}
\label{tab:spearman-npp}
\begin{tabular}{@{}lrl@{}}
\toprule
Segmento & $\rho$ de Spearman & Significância \\
\midrule
Pecuária             & $+0{,}52$ & $p < 0{,}001$ \\
Conservação          & $+0{,}46$ & $p < 0{,}001$ \\
Cana-de-açúcar       & $+0{,}45$ & $p < 0{,}001$ \\
Geração fotovoltaica & $+0{,}40$ & $p < 0{,}001$ \\
Soja                 & $+0{,}31$ & $p < 0{,}001$ \\
Outras culturas anuais & $+0{,}11$ & n.s. ($p = 0{,}09$) \\
Piscicultura         & $-0{,}32$ & $p < 0{,}001$ \\
\bottomrule
\end{tabular}
\end{table}

O \emph{proxy} de NPP anual mede a produtividade primária da vegetação, não o rendimento
agronômico, de modo que os segmentos de biomassa em pé (pecuária, conservação) pontuam mais
alto. Por isso ele é relatado apenas como \emph{contexto} de produtividade; as métricas de
discriminação por presença (Tabela~\ref{tab:presence-auc}) são os validadores primários das
culturas.

\chapter{Figuras de diagnóstico}\label{anx:figures}


\begin{figure}[H]
\centering
\includegraphics[width=0.85\textwidth]{fig_4_2.png}
\caption{Classes de viabilidade da FAO (N / S3 / S2 / S1) para segmentos representativos
(soja e conservação).}
\label{fig:fao-classes}
\end{figure}

\begin{figure}[H]
\centering
\includegraphics[width=0.7\textwidth]{fig_4_6.png}
\caption{Realocação de melhor cultura --- células em que muda a cultura mais viável entre
soja, cana-de-açúcar e outras culturas anuais --- sob SSP5-8.5, 2051--2070.}
\label{fig:best-crop}
\end{figure}

\begin{figure}[H]
\centering
\includegraphics[width=0.7\textwidth]{fig_4_8.png}
\caption{Subutilização potencial \emph{versus} realizado para a soja --- células viáveis ou
superiores mas não atualmente sob soja.}
\label{fig:underused-soy}
\end{figure}

\begin{figure}[H]
\centering
\includegraphics[width=\textwidth]{fig_4_13.png}
\caption{Viabilidade média municipal para os sete segmentos --- o equivalente municipal da Figura~\ref{fig:atlas-present},
agregado por redução zonal sobre os municípios; complemento da Figura~\ref{fig:municipal}, que só detalha soja e outras culturas.}
\label{fig:all-segment-suit}
\end{figure}

\section{Seleção do número de zonas (k)}\label{sec:k-selection-annex}

\begin{figure}[H]
\centering
\includegraphics[width=0.9\textwidth]{fig_4_10.png}
\caption{Índices de validade de agrupamento (silhueta, Davies--Bouldin, estatística de
\emph{gap}) e estabilidade em subamostras (índice de Rand ajustado) ao longo de
$k = 2\dots20$ no espaço de componentes principais descorrelacionado.}
\label{fig:cluster-validity}
\end{figure}

A silhueta é quase plana (0,161--0,187) ao longo de todo o intervalo; o Davies--Bouldin cai até
$k=7$ e fica praticamente plano de $k=7$ a $20$ (1,50--1,57, mínimo em $k=19$); e o \emph{gap}
cresce quase monotonicamente, estabilizando-se só em $k=14$. Nenhum índice isolado decide.
$k = 4$ (linha tracejada) é o máximo de silhueta de todo o \emph{sweep} (0,210, contra 0,206
em $k=5$) \emph{e} o máximo de estabilidade sob subamostragem (ARI $=0{,}988$), com o menor
desvio-padrão entre réplicas de todo o intervalo ($\pm$0,004; a título de comparação, $\pm$0,096
em $k=5$ e $\pm$0,124 em $k=6$). O Davies--Bouldin prefere $k=5$ (1,463 contra 1,574), a única
discordância entre os índices. O painel inteiro varia pouco: a silhueta fica entre 0,158 e 0,210
em $k = 2,\dots,20$, o que indica ausência de estrutura de agrupamento forte em qualquer $k$ e
é a razão de a escolha se apoiar na reprodutibilidade da partição, e não num ótimo de validade
(\S\ref{sec:zoning}). O painel cobre $k$ até 20, confirmando que a escolha não é um artefato de
faixa truncada.

\section{Decomposição por fator}\label{sec:factor-variance-annex}

\begin{sidewaysfigure}[htp]
\centering
\includegraphics[width=0.95\textheight]{fig_4_11.png}
\caption{Parcela de cada fator na variância espacial da viabilidade de cada segmento, por
decomposição exata de covariância entre cada fator ponderado e o agregado (as parcelas somam
exatamente~1 por construção). Eixo horizontal em escala simétrica-log; a cor indica o grupo
temático do fator (clima, laranja; terreno+solo, verde; outros --- água, sítio, acesso,
extras ---, azul); contorno preto grosso marca o fator dominante de cada segmento; hachura
indica contribuição negativa.}
\label{fig:factor-variance}
\end{sidewaysfigure}

A decomposição é feita sobre log-pertinências para os seis segmentos de média geométrica e sobre
pertinências brutas para a conservação, de média aritmética; o eixo horizontal usa uma escala
simétrica-log (limiar linear~$=0{,}005$) para manter visíveis as parcelas próximas de zero sem
exagerar sua magnitude real, e a contribuição negativa indicada por hachura é um efeito residual
esperado de uma decomposição linear de covariância entre fatores correlacionados, não um erro.

Terreno e solo dominam as culturas, mas o fator vencedor não é sempre o mesmo: declive responde
por 0,78 da variância da soja (argila, 0,22) e por $\approx$1,0 de outras culturas anuais, ao
passo que para a cana-de-açúcar é a argila que domina (0,69, contra 0,31 do declive). Para soja e
outras culturas anuais, apenas argila (e declive) de fato discriminam --- pH e carbono orgânico
do solo também saturam (parcela $\approx$0), de modo que \emph{solo} isoladamente é um
discriminador espacial mais estreito do que \emph{terreno e solo} em conjunto. A pecuária é
governada pela umidade do solo (0,58) e pelo carbono orgânico (0,26); a piscicultura reparte-se
entre distância à água sazonal (0,44), declive (0,41) e densidade de drenagem (0,15); a
conservação, o segmento mais equilibrado dos sete, divide-se em três fatores de peso
comparável --- conectividade a áreas protegidas (0,31), rugosidade do terreno (0,23) e fração
de vegetação nativa (0,22); a geração fotovoltaica, pelo índice de claridade (0,39) e pela
radiação solar (0,36). Essa concentração não decorre de viés de ponderação AHP, mas de
saturação geográfica genuína das pertinências climáticas, confirmada de forma independente pelo
mapa de magnitude espacial da Figura~\ref{fig:scale-mismatch} (argumento completo em
\S\ref{sec:atlas-present}).

\section{Descompasso de escala clima--terreno}\label{sec:scale-mismatch-annex}

\begin{figure}[H]
\centering
\includegraphics[width=0.95\textwidth]{fig_4_12.png}
\caption{Variabilidade espacial local (desvio-padrão em janela móvel, bandas padronizadas)
do clima \emph{versus} terreno+solo, numa janela de 4,75~km --- aproximadamente o tamanho da
célula nativa do TerraClimate ($\tfrac{1}{24}^\circ \approx 4{,}64$~km).}
\label{fig:scale-mismatch}
\end{figure}

\begin{table}[H]
\centering
\caption{Desvio-padrão local médio (bandas padronizadas) do clima \emph{versus}
terreno+solo, por raio de janela móvel, e razão terreno+solo/clima correspondente.}
\label{tab:scale-mismatch}
\begin{tabular}{@{}rrrrr@{}}
\toprule
Raio (m) & $\sigma$ clima & $\sigma$ terreno+solo & Razão & $\log_2$(razão) \\
\midrule
250 & 0,001 & 0,437 & 320 & 8,33 \\
750 & 0,003 & 0,655 & 198 & 7,63 \\
4.750 & 0,021 & 0,900 & 42 & 5,40 \\
\bottomrule
\end{tabular}
\end{table}

O painel do clima (Figura~\ref{fig:scale-mismatch}) é quase uniformemente plano (variação local
desprezível em toda a extensão do território); o painel de terreno+solo mostra estrutura
espacial genuína, mais acentuada no norte/nordeste (transição de relevo do planalto); o painel
da razão log$_2$ é quase uniformemente positivo, confirmando que terreno+solo domina a estrutura
espacial fina mesmo na própria escala nativa do clima. O mapa é construído a partir da mesma
amostra limitada por orçamento computacional (4.475 pontos por raio, \S\ref{anx:data}) que
alimenta a Tabela~\ref{tab:scale-mismatch}, e não de uma composição rasterizada completa do
território. Nessa tabela, o raio de 4.750~m ancora na célula nativa do TerraClimate
($\approx$4,64~km); 250~m é a célula nativa da pilha de atributos; e 750~m reutiliza o raio já
usado para a rugosidade do terreno na conservação (\S\ref{anx:derivations}).
